\subsection{ソフトウェア保守の基礎}
この節では、ソフトウェア保守の役割と範囲を理解するための概念と用語を整理する。特に重要なのは、保守を種類別に分類して理解することである。分類ができると、費用、優先順位、必要な人員、テスト範囲、リリース判断を説明しやすくなる。\par
\subsubsection{定義と用語}
ソフトウェア保守の目的は、既存ソフトウェアを変更しながら、その完全性を保つことである。完全性とは、単に動くことではなく、要求、設計、コード、テスト、文書、構成、運用上の整合が保たれている状態を意味する。\par
SWEBOK では、ソフトウェア保守を「運用中のソフトウェアに費用対効果の高い支援を提供するために必要な活動の総体」として扱う。したがって、保守には納品後の修正だけでなく、納品前の保守計画、保守性の確保、移行支援、訓練、ヘルプデスク連携も含まれる。\par
\begin{definitionbox}ソフトウェア保守とは、運用中のソフトウェアを支援し、変更し、品質と価値を保つための活動である。変更しながら壊さないこと、すなわち「変化」と「完全性」の両立が中心課題である。\end{definitionbox}
\subsubsection{ソフトウェア保守の性質}
\term{保守}{Maintenance}は、開発から運用、最終的な廃止まで、ソフトウェアのライフサイクル全体を支える。運用中のソフトウェアは、容量、継続性、可用性の観点で監視される。変更要求や問題報告は記録され、追跡され、影響が分析される。その後、コードや文書などの成果物が修正され、テストされ、新しい版が運用へリリースされる。\par
保守担当者は、開発者や運用者が持つ知識から学ぶ必要がある。保守担当者が早期に開発へ関与すれば、保守費用と手戻りを減らしやすい。反対に、開発者が離れた後に保守担当者が参加すると、コード、テスト、文書、設計意図を理解する負担が大きくなる。\par
\subsubsection{ソフトウェア保守が必要な理由}
\term{ソフトウェア保守}{Software Maintenance}は、ソフトウェアが寿命の間ずっと利用者要求を満たし続けるために必要である。開発ライフサイクルがウォーターフォール型でもアジャイル型でも、ソフトウェアは変更される。変更の理由は、欠陥修正だけではない。\par
\begin{itemize}
\item 欠陥や潜在欠陥を修正するため。
\item 運用中ソフトウェアの設計、性能、信頼性を改善するため。
\item 新しい機能や利用者要求を実装するため。
\item 利用者が機能を理解し、使い続けられるように支援するため。
\item 外部システム、基盤、ハードウェア、クラウド、法規制の変化へ適応するため。
\item セキュリティ脅威を予防し、脆弱性を修正するため。
\item 技術的陳腐化に対応し、老朽化した要素を更新するため。
\item 不要になったソフトウェアを安全に廃止するため。
\end{itemize}
\subsubsection{保守費用が大きい理由}
ソフトウェアの総費用の中で、保守は大きな割合を占める。一般に、誤りの修正費用はライフサイクルの後半ほど大きくなる。また、保守は長期間続くため、単発の開発費用よりも見えにくいが、累積費用は大きい。\par
保守は「欠陥修正が中心」と誤解されやすい。しかし、調査では保守の多くが機能強化や環境適応に使われるとされる。修正と強化を同じ報告カテゴリにまとめると、欠陥修正が実際以上に高く見える。保守分類を分けて測定することは、費用構造を理解するために重要である。\par
\begin{pointbox}{費用を見る観点}保守費用は、運用環境と組織環境の両方に左右される。運用環境にはハードウェア、OS、ミドルウェア、クラウド基盤が含まれる。組織環境には方針、競争、プロセス、製品、人員が含まれる。\end{pointbox}
\subsubsection{ソフトウェアの進化}
ソフトウェア保守は、ソフトウェア進化を支える活動である。Lehman らの研究は、長期運用されるソフトウェアが変化し続けること、そして変化を管理しなければ複雑さや品質低下が進むことを示した。\par
\par\smallskip\noindent\refstepcounter{table}\label{tab:ch07-02}表 \thetable: ソフトウェアの進化における進化の法則・初学者向けの意味\par\smallskip
表 \thetable は、ソフトウェアの進化における進化の法則・初学者向けの意味を比較・確認しやすい形で整理したものである。\par\smallskip
\begin{longtable}{p{0.30\linewidth}p{0.64\linewidth}}
\toprule
進化の法則 & 初学者向けの意味 \\
\midrule
\endhead
継続的変化 & ソフトウェアは利用環境に合わせて継続的に適応しなければ、満足度が下がる。 \\
複雑さの増大 & 変更を重ねると、複雑さは増える。複雑さを減らす作業をしなければ保守は難しくなる。 \\
自己調整 & 進化の過程は、製品やプロセスの測定値を通じて一定の傾向を持つ。 \\
作業率の不変性 & 長期的には、進化するソフトウェアに投入できる有効作業量には一定の制約がある。 \\
親熟度の保存 & 関係者がソフトウェアの内容と振る舞いを理解し続ける必要がある。成長が急すぎると理解が追いつかない。 \\
継続的成長 & 利用者満足を保つため、機能内容は継続的に増える傾向がある。 \\
品質低下 & 運用環境の変化へ厳密に適応しなければ、品質は低下して見える。 \\
フィードバックシステム & ソフトウェア進化は複数の利害関係者と情報の循環から成るため、フィードバックとして扱う必要がある。 \\
\bottomrule
\end{longtable}
\par\smallskip
\begin{center}
\centering
\IfFileExists{assets/generated_figures/ch07/fig-ch07-02.png}{\includegraphics[width=0.96\linewidth]{assets/generated_figures/ch07/fig-ch07-02.png}}{\fbox{\parbox[c][48mm][c]{0.9\linewidth}{\centering 画像生成待ち\\ソフトウェア進化の循環図}}}
\par\smallskip
\refstepcounter{figure}\label{fig:ch07-02}図 \thefigure: ソフトウェア進化の循環図
\end{center}
図 \thefigure は、ソフトウェア進化の循環図を視覚的に整理したものである。
\par\smallskip
\subsubsection{ソフトウェア保守の分類}
ソフトウェア保守は、変更要求を分類するために複数の種類へ分けられる。SWEBOK では、是正、予防、適応、追加、完全化に加えて、緊急保守を説明する。\par
\par\smallskip\noindent\refstepcounter{table}\label{tab:ch07-03}表 \thetable: ソフトウェア保守の分類における分類・説明・例\par\smallskip
表 \thetable は、ソフトウェア保守の分類における分類・説明・例を比較・確認しやすい形で整理したものである。\par\smallskip
\begin{longtable}{p{0.22\linewidth}p{0.40\linewidth}p{0.30\linewidth}}
\toprule
分類 & 説明 & 例 \\
\midrule
\endhead
是正保守 & 納品後に見つかった問題を修正する反応的な変更である。 & 本番障害の原因となった計算誤りを修正する。 \\
予防保守 & 本番で問題が起きる前に潜在欠陥を修正する変更である。 & 将来エラーになり得る古い日付処理を直す。 \\
適応保守 & 環境変化の中でソフトウェアを使い続けられるようにする変更である。 & OS、データベース、外部APIの変更に対応する。 \\
追加保守 & 利用を拡張するために、比較的大きな機能や特徴を追加する変更である。 & 既存サービスに新しい通知機能を追加する。 \\
完全化保守 & 利用者向け改善、文書改善、性能や保守性などの品質改善を行う変更である。 & 画面操作を簡単にする、コードを読みやすくする。 \\
緊急保守 & 是正保守までの間、一時的にシステムを稼働させ続けるための予定外変更である。 & 本番停止を避けるため、一時的な回避策を適用する。 \\
\bottomrule
\end{longtable}
\begin{pointbox}{分類の使い方}保守分類は、会計上の分類ではなく、意思決定の道具である。どの保守が多いかを測れば、品質問題、環境変化、利用者要求、技術的負債のどこに費用が使われているかが見える。\end{pointbox}
\par\smallskip
\begin{center}
\centering
\IfFileExists{assets/generated_figures/ch07/fig-ch07-03.png}{\includegraphics[width=0.96\linewidth]{assets/generated_figures/ch07/fig-ch07-03.png}}{\fbox{\parbox[c][48mm][c]{0.9\linewidth}{\centering 画像生成待ち\\ソフトウェア保守分類図}}}
\par\smallskip
\refstepcounter{figure}\label{fig:ch07-03}図 \thefigure: ソフトウェア保守分類図
\end{center}
図 \thefigure は、ソフトウェア保守分類図を視覚的に整理したものである。
\par\smallskip
